\documentclass[10pt,a4paper]{amsart}
\usepackage{amsmath,amssymb,amsthm,graphicx}
\usepackage{verbatim}
\newenvironment{reference}{\comment}{\endcomment}
%\newenvironment{reference}{}{}
\newenvironment{slogan}{\comment}{\endcomment}
\newenvironment{history}{\comment}{\endcomment}

% For commutative diagrams we use Xy-pic
\usepackage[all]{xy}

% We use 2cell for 2-commutative diagrams.
\xyoption{2cell}
% Theorem environments.
%
\theoremstyle{plain}
\newtheorem{theorem}[subsection]{Theorem}
\newtheorem{proposition}[subsection]{Proposition}
\newtheorem{lemma}[subsection]{Lemma}

\theoremstyle{definition}
\newtheorem{definition}[subsection]{Definition}
\newtheorem{example}[subsection]{Example}
\newtheorem{exercise}[subsection]{Exercise}
\newtheorem{situation}[subsection]{Situation}

\theoremstyle{remark}
\newtheorem{remark}[subsection]{Remark}
\newtheorem{remarks}[subsection]{Remarks}

\numberwithin{equation}{subsection}

% Macros
%
\def\lim{\mathop{\mathrm{lim}}\nolimits}
\def\colim{\mathop{\mathrm{colim}}\nolimits}
\def\Spec{\mathop{\mathrm{Spec}}}
\def\Hom{\mathop{\mathrm{Hom}}\nolimits}
\def\Ext{\mathop{\mathrm{Ext}}\nolimits}
\def\SheafHom{\mathop{\mathcal{H}\!\mathit{om}}\nolimits}
\def\SheafExt{\mathop{\mathcal{E}\!\mathit{xt}}\nolimits}
\def\Sch{\mathit{Sch}}
\def\Mor{\mathop{\mathrm{Mor}}\nolimits}
\def\Ob{\mathop{\mathrm{Ob}}\nolimits}
\def\Sh{\mathop{\mathit{Sh}}\nolimits}
\def\NL{\mathop{N\!L}\nolimits}
\def\CH{\mathop{\mathrm{CH}}\nolimits}
\def\proetale{{pro\text{-}\acute{e}tale}}
\def\etale{{\acute{e}tale}}
\def\QCoh{\mathit{QCoh}}
\def\Ker{\mathop{\mathrm{Ker}}}
\def\Im{\mathop{\mathrm{Im}}}
\def\Coker{\mathop{\mathrm{Coker}}}
\def\Coim{\mathop{\mathrm{Coim}}}

% Boxtimes
%
\DeclareMathSymbol{\boxtimes}{\mathbin}{AMSa}{"02}

%
% Macros for moduli stacks/spaces
%
\def\QCohstack{\mathcal{QC}\!\mathit{oh}}
\def\Cohstack{\mathcal{C}\!\mathit{oh}}
\def\Spacesstack{\mathcal{S}\!\mathit{paces}}
\def\Quotfunctor{\mathrm{Quot}}
\def\Hilbfunctor{\mathrm{Hilb}}
\def\Curvesstack{\mathcal{C}\!\mathit{urves}}
\def\Polarizedstack{\mathcal{P}\!\mathit{olarized}}
\def\Complexesstack{\mathcal{C}\!\mathit{omplexes}}
% \Pic is the operator that assigns to X its picard group, usage \Pic(X)
% \Picardstack_{X/B} denotes the Picard stack of X over B
% \Picardfunctor_{X/B} denotes the Picard functor of X over B
\def\Pic{\mathop{\mathrm{Pic}}\nolimits}
\def\Picardstack{\mathcal{P}\!\mathit{ic}}
\def\Picardfunctor{\mathrm{Pic}}
\def\Deformationcategory{\mathcal{D}\!\mathit{ef}}
\usepackage{cleveref}
\crefname{originalproposition}{original proposition}{original propositions}
\Crefname{originalproposition}{Original proposition}{Original propositions}
\makeatletter
\newlabel{cohomology-proposition-vanishing-Noetherian}{{02UZ}{source}{}{sourcepointer-1}{}}
\newlabel{cohomology-proposition-vanishing-Noetherian@cref}{{[originalproposition][0][]02UZ}{[1][0][]source}{}{}{}}
\makeatother
\crefname{originallemma}{original lemma}{original lemmas}
\Crefname{originallemma}{Original lemma}{Original lemmas}
\makeatletter
\newlabel{more-algebra-lemma-lift-perfect-from-residue-field}{{0BCD}{source}{}{sourcepointer-2}{}}
\newlabel{more-algebra-lemma-lift-perfect-from-residue-field@cref}{{[originallemma][0][]0BCD}{[1][0][]source}{}{}{}}
\makeatother
\crefname{originallemma}{original lemma}{original lemmas}
\Crefname{originallemma}{Original lemma}{Original lemmas}
\makeatletter
\newlabel{perfect-lemma-flat-proper-perfect-direct-image-general}{{0B91}{source}{}{sourcepointer-3}{}}
\newlabel{perfect-lemma-flat-proper-perfect-direct-image-general@cref}{{[originallemma][0][]0B91}{[1][0][]source}{}{}{}}
\makeatother
\crefname{originallemma}{original lemma}{original lemmas}
\Crefname{originallemma}{Original lemma}{Original lemmas}
\makeatletter
\newlabel{perfect-lemma-proper-flat-geom-red-connected}{{0E0L}{source}{}{sourcepointer-4}{}}
\newlabel{perfect-lemma-proper-flat-geom-red-connected@cref}{{[originallemma][0][]0E0L}{[1][0][]source}{}{}{}}
\makeatother
\crefname{originallemma}{original lemma}{original lemmas}
\Crefname{originallemma}{Original lemma}{Original lemmas}
\makeatletter
\newlabel{perfect-lemma-proper-flat-h0}{{0E62}{source}{}{sourcepointer-5}{}}
\newlabel{perfect-lemma-proper-flat-h0@cref}{{[originallemma][0][]0E62}{[1][0][]source}{}{}{}}
\makeatother
\title{Cohomology of the structure sheaf in a nodal family with connected geometric fibres: pushforwards, local freeness and base change}
\author{The Stacks Project Authors}
\date{stacks-project:curves.tex:local-snapshot-20261012; GNU FDL 1.2 or later}
\begin{document}
\maketitle
\noindent\textit{Editorial note.} All delivered mathematical text is the author's own, copied from the pinned Stacks Project chapter file curves.tex; no statement or proof step was rewritten, reordered or removed.\par\medskip
\noindent\textit{Editorial note.} The counted proof is the author's complete argument: the pushforward identification, the perfect decomposition with tor amplitude, the fibre cohomology computation with its vanishing, and the finite-local-freeness conclusion with base-change compatibility are all retained verbatim.\par\medskip
\noindent\textit{Editorial note.} The label lemma-genus-in-nodal-family-of-curves names the counted statement, and the proof environment immediately following it in the source is delivered verbatim as its proof.\par\medskip
\noindent\textit{Editorial note.} The counted statement is carried by the author's own lemma environment in the chapter source.\par\medskip
\noindent\textit{Allowed standard prerequisites (editorial).}\begin{itemize}
\item Pushforward for proper flat morphisms with geometrically reduced connected fibres. Derived Categories of Schemes, Lemma perfect-lemma-proper-flat-geom-red-connected. Gives part (a): f\_*O\_X = O\_S compatible with base change.
\item Perfect decomposition of the total pushforward. Derived Categories of Schemes, Lemmas perfect-lemma-proper-flat-h0 and perfect-lemma-flat-proper-perfect-direct-image-general. Writes Rf\_*O\_X = O\_S (+) P locally with P perfect of tor amplitude in [1, infinity), and gives base-change compatibility of Rf\_*O\_X.
\item Vanishing and perfect lifting. Cohomology, Proposition cohomology-proposition-vanishing-Noetherian; More on Algebra, Lemma more-algebra-lemma-lift-perfect-from-residue-field. Kills the fibre cohomology in degrees at least 2 on the one-dimensional Noetherian fibres and promotes H\textasciicircum{}1(P) to a finite locally free module from its residue-field fibres.
\end{itemize}\medskip
\ifdefined\hypertarget\hypertarget{sourcepointer-1}{}\fi
\noindent\textit{Source reference.} Original proposition 02UZ: sources/stacks-chow/tags.txt, source lines 3750--3750. Prerequisite pointer invoked at proof line(s) 5268; it names a standard fact of a sibling chapter that the counted argument uses, and every step of the counted proof built on it is retained verbatim in the delivered body.\par\medskip
\ifdefined\hypertarget\hypertarget{sourcepointer-2}{}\fi
\noindent\textit{Source reference.} Original lemma 0BCD: sources/stacks-chow/tags.txt, source lines 14276--14276. Prerequisite pointer invoked at proof line(s) 5271; it names a standard fact of a sibling chapter that the counted argument uses, and every step of the counted proof built on it is retained verbatim in the delivered body.\par\medskip
\ifdefined\hypertarget\hypertarget{sourcepointer-3}{}\fi
\noindent\textit{Source reference.} Original lemma 0B91: sources/stacks-chow/tags.txt, source lines 14159--14159. Prerequisite pointer invoked at proof line(s) 5259; it names a standard fact of a sibling chapter that the counted argument uses, and every step of the counted proof built on it is retained verbatim in the delivered body.\par\medskip
\ifdefined\hypertarget\hypertarget{sourcepointer-4}{}\fi
\noindent\textit{Source reference.} Original lemma 0E0L: sources/stacks-chow/tags.txt, source lines 17541--17541. Prerequisite pointer invoked at proof line(s) 5251; it names a standard fact of a sibling chapter that the counted argument uses, and every step of the counted proof built on it is retained verbatim in the delivered body.\par\medskip
\ifdefined\hypertarget\hypertarget{sourcepointer-5}{}\fi
\noindent\textit{Source reference.} Original lemma 0E62: sources/stacks-chow/tags.txt, source lines 17732--17732. Prerequisite pointer invoked at proof line(s) 5253; it names a standard fact of a sibling chapter that the counted argument uses, and every step of the counted proof built on it is retained verbatim in the delivered body.\par\medskip
\begin{lemma}
\label{lemma-genus-in-nodal-family-of-curves}
Let $f : X \to S$ be a morphism of schemes. Assume
\begin{enumerate}
\item $f$ is proper,
\item $f$ is at-worst-nodal of relative dimension $1$, and
\item the geometric fibres of $f$ are connected.
\end{enumerate}
Then (a) $f_*\mathcal{O}_X = \mathcal{O}_S$ and this holds after
any base change, (b) $R^1f_*\mathcal{O}_X$ is a finite locally free
$\mathcal{O}_S$-module whose formation commutes with any base change,
and (c) $R^qf_*\mathcal{O}_X = 0$ for $q \geq 2$.
\end{lemma}

\begin{proof}
Part (a) follows from
Derived Categories of Schemes, Lemma
\ref{perfect-lemma-proper-flat-geom-red-connected}.
By Derived Categories of Schemes, Lemma
\ref{perfect-lemma-proper-flat-h0} locally on $S$
we can write $Rf_*\mathcal{O}_X = \mathcal{O}_S \oplus P$
where $P$ is perfect of tor amplitude in $[1, \infty)$.
Recall that formation of $Rf_*\mathcal{O}_X$ commutes
with arbitrary base change
(Derived Categories of Schemes, Lemma
\ref{perfect-lemma-flat-proper-perfect-direct-image-general}).
Thus for $s \in S$ we have
$$
H^i(P \otimes_{\mathcal{O}_S}^\mathbf{L} \kappa(s)) =
H^i(X_s, \mathcal{O}_{X_s})
\text{ for }i \geq 1
$$
This is zero unless $i = 1$ since $X_s$ is a $1$-dimensional
Noetherian scheme, see
Cohomology, Proposition \ref{cohomology-proposition-vanishing-Noetherian}.
Then $P = H^1(P)[-1]$ and $H^1(P)$ is finite locally free
for example by More on Algebra, Lemma
\ref{more-algebra-lemma-lift-perfect-from-residue-field}.
Since everything is compatible with base change we conclude.
\end{proof}

\end{document}
